%% 2i2d-descente-charges — le cheminement des charges dans un bâtiment :
%% charges → porteur horizontal (flexion) → porteurs verticaux (compression)
%% → fondations → sol (réaction répartie).
\documentclass[tikz,border=4pt]{standalone}
\usepackage{liaisons}
\begin{document}
\begin{tikzpicture}[x=1cm,y=1cm,
  charge/.style={-{Stealth[length=4pt]}, line width=0.8pt, solideA},
  reaction/.style={-{Stealth[length=4pt]}, line width=0.8pt, solideF},
  petit/.style={font=\scriptsize, inner sep=1.5pt}]

\newcommand\batihoriz[3]{%
  \draw[line width=1pt, solideE] (#1,#3) -- (#2,#3);
  \begin{scope}
    \clip (#1,#3-0.26) rectangle (#2,#3);
    \foreach \hx in {-16,-15.76,...,16} { \draw[line width=0.5pt, solideE!75] (\hx,#3) -- (\hx-0.3,#3-0.3); }
  \end{scope}}

%% --- 1. les charges qui arrivent sur la dalle ---------------------------------
\foreach \x in {1.0,1.62,2.24,2.85,3.47,4.09,4.7} { \draw[charge] (\x,4.45) -- (\x,3.92); }
\node[petit, anchor=west, text=solideA, align=left] at (5.35,4.2)
     {charges : poids propre,\\exploitation, neige};

%% --- 2. porteur horizontal : la dalle, qui fléchit ------------------------------
\draw[line width=0.7pt, draw=solideB, fill=solideB!18] (0.7,3.5) rectangle (5.0,3.85);
\draw[line width=0.7pt, solideB!60] (0.7,3.44) .. controls (2.1,3.2) and (3.6,3.2) .. (5.0,3.44);
\node[petit, anchor=west, text=solideB, align=left] at (5.35,3.5)
     {porteur horizontal\\$\rightarrow$ \textbf{flexion}};

%% --- 3. porteurs verticaux : les poteaux, comprimés -------------------------------
\foreach \x in {1.05,4.2} {
  \draw[line width=0.7pt, draw=solideB, fill=solideB!18] (\x,1.3) rectangle (\x+0.45,3.5);
  \draw[charge] (\x+0.225,3.35) -- (\x+0.225,2.75);
  \draw[charge] (\x+0.225,1.45) -- (\x+0.225,2.05);}
\node[petit, anchor=west, text=solideB, align=left] at (5.35,2.4)
     {porteur vertical\\$\rightarrow$ \textbf{compression}};

%% --- 4. fondations : les semelles ---------------------------------------------------
\draw[line width=0.7pt, draw=solideE, fill=solideE!20, line join=round]
  (0.7,0.9) -- (1.85,0.9) -- (1.6,1.3) -- (0.95,1.3) -- cycle;
\draw[line width=0.7pt, draw=solideE, fill=solideE!20, line join=round]
  (3.85,0.9) -- (5.0,0.9) -- (4.75,1.3) -- (4.1,1.3) -- cycle;
\node[petit, anchor=west, text=solideE] at (5.35,1.1) {fondations};

%% --- 5. le sol et sa réaction ----------------------------------------------------------
\batihoriz{0.45}{5.25}{0.9}
\foreach \x in {0.85,1.15,1.45,1.75,3.95,4.25,4.55,4.85} { \draw[reaction] (\x,0.28) -- (\x,0.82); }
\node[petit, anchor=west, text=solideF] at (5.35,0.5) {réaction du sol};

%% --- la grande flèche « descente de charges » -----------------------------------------
\draw[-{Stealth[length=7pt]}, line width=2.2pt, solideD] (0.2,3.9) -- (0.2,0.98);
\node[petit, anchor=south, text=solideD, rotate=90] at (0.06,1.55) {descente de charges};
\end{tikzpicture}
\end{document}
